\section{Cross-Sectional Wasserstein Dispersion and Spatial Attenuation}
\label{sec:dispersion}

The dispersion analysis is the only part of the paper that uses a free-centre
problem.
Unlike the fixed-target reconstruction in \Cref{sec:interaction-operator}, the
unrestricted Wasserstein barycentre problem below selects a centre distribution
$Q$ that summarizes heterogeneity across firms.

How much characteristic-implied exposure heterogeneity survives peer adjustment?
The conditional answer is economically simple: characteristic
dispersion places a floor on stand-alone exposure dispersion, and
adjustment through an admissible peer field cannot erase more than a
coefficient-dependent share of that floor.
This is a supporting closure result for the field constructed in
\Cref{sec:interaction-operator}, not a second field construction or
an empirical calibration of the carrier restrictions introduced below.

Fix cross-sectional weights $q=(q_1,\ldots,q_N)$ in the unit simplex.
Assume that the laws considered below have finite second moments.
The simplex restriction makes the weights nonnegative with unit sum,
so they determine each firm's contribution without changing the scale
of the aggregate comparison.
Finite second moments make the quadratic transport costs below well defined.
Economically, $\mathcal D_q(C_1,\dots,C_N)$ is the least weighted
separation compatible with all observed firm-level characteristic laws.
It is therefore a conservative measure of cross-sectional information
heterogeneity rather than the dispersion generated by one selected matching.
Formally, for characteristic laws on $(\Omega,d_\Omega)$ and
$(X_1,\dots,X_N)\sim\gamma$, define
\[
  \mathcal D_q(C_1,\dots,C_N)
  =\inf_{\gamma\in\Pi}
  \mathbb E_\gamma\Big[\sum_{i<j}q_{i}q_{j}\,d_\Omega(X_i,X_j)^2\Big],
\]
with $\Pi$ the joint laws on $\Omega^N$ whose $i$th marginal is $C_i$.
The infimum asks how close the laws could be under their most
favourable common coupling.
Its units are squared embedding distance, and its value depends on
the chosen ground metric and weights.

The Hilbert-space restriction begins to matter for the centre
representation: for any laws $R_1,\ldots,R_N$ on a Hilbert space, the
same functional satisfies
$\mathcal D_q(R_1,\ldots,R_N)=\inf_Q\sum_i q_i\,W_2^2(R_i,Q)$
\citep{agueh_carlier_barycenters_2011,gawronsky_distance_mpt_2026}.
Here $Q$ is free and summarizes cross-sectional dispersion.
Its minimizer is the unrestricted Wasserstein barycentre of the firm laws under
weights $q$.

For two firms, the aggregate functional reduces to the familiar
pairwise Wasserstein distance.

\begin{theorem}[Pairwise nesting]\label{thm:pairwise-nesting}
  For $N=2$,
  \begin{equation}\label{eq:pairwise-nesting}
    \mathcal D_{(q_1,q_2)}(C_1,C_2)=q_1q_2\,W_2^2(C_1,C_2).
  \end{equation}
\end{theorem}
For two marginals, $\Pi$ is exactly the set of couplings of $C_1,C_2$.
The objective is therefore $q_1q_2$ times the quadratic transport
cost, so taking the infimum gives Equation~\eqref{eq:pairwise-nesting}.

Thus the geometric object nests the squared Wasserstein term used to
derive the pairwise covariance envelope in \citet{gawronsky_continuous_2026}.
The equality establishes geometric nesting only; it does not import
that paper's covariance or coupling assumptions.

The Wasserstein dispersion measure is in embedding units, whereas
stand-alone exposures are in factor-loading units, so an explicit
transfer restriction is needed before the two can be compared.
Write $Z_i=\Gamma^{1/2}\xi_i$ and $P_i=\mathcal L(Z_i)$ for the
stand-alone exposure laws in covariance coordinates.
The transformation by $\Gamma^{1/2}$ places exposure differences in
the risk coordinates used by the dispersion certificate.
Let $t:\Omega\to\mathcal H$ be a common measurable carrier in those coordinates.
Commonness supplies one benchmark across firms, and measurability
makes the carrier image of each information draw a valid random exposure.
For some $L>0$, impose the noncollapse condition at the point where
characteristic distance must become exposure distance:
\[
  \norm{t(x)-t(y)}\geq L^{-1}d_\Omega(x,y)
  \quad\text{for every }x,y\in\Omega.
\]
This lower-distance restriction prevents the common carrier from
erasing separation between distinct information states.
No upper-distance bound enters the result.
For each firm, next allow synchronous slack $\tau_i\geq0$:
\[
  \norm{\Gamma^{1/2}T(X_i,U_i)-t(X_i)}\leq\tau_i
  \quad\text{almost surely}.
\]
The same draw $X_i$ appears on both sides because the bound must
couple firm $i$'s actual stand-alone exposure to the carrier image of
that information state.
The slack permits firm-specific transmission noise of at most
$\tau_i$ in risk-coordinate norm.
Together, noncollapse and synchronous slack state the entire
characteristic-to-exposure restriction used below.
They do not restrict $W$ or the peer-adjustment mechanism.

Writing $\tau_q=(\sum_i q_i\tau_i^2)^{1/2}$ for weighted
root-mean-square slack and $[a]_+=\max\{a,0\}$, we maintain the
following characteristic-to-exposure transfer inequality:
\begin{equation}\label{eq:imported-transfer}
  \sqrt{\mathcal D_q(P_1,\dots,P_N)}
  \geq
  \left[L^{-1}\sqrt{\mathcal D_q(C_1,\dots,C_N)}-\tau_q\right]_+.
\end{equation}
The root-mean-square aggregation matches the cross-sectional weights,
while the positive part sets the lower floor to zero when
transmission slack absorbs the carrier-adjusted separation.
This restriction is stated here in full;
\citet{gawronsky_distance_mpt_2026} studies its portfolio implications.

Observable separation is informative whenever its carrier-adjusted
magnitude exceeds the root-mean-square transmission slack.
At zero slack \eqref{eq:imported-transfer} becomes $\mathcal
D_q(P_1,\ldots,P_N)\geq L^{-2}\mathcal D_q(C_1,\ldots,C_N)$.
The carrier, $L$, and $\tau_i$ are maintained inputs rather than
estimated objects, so \eqref{eq:imported-transfer} is a conditional
restriction and is not numerically calibrated here.

\subsection{Spatial attenuation under peer adjustment}

The transfer inequality stops at stand-alone exposure dispersion;
attenuation enters only to ask what the spatial closure does to that floor.
The reduced form in \Cref{thm:sar-reduced-form} maps stand-alone
exposure $\xi$ into peer-adjusted exposure $B$ through the normalized
spatial multiplier:

\begin{equation}\label{eq:sar-multiplier}
  S_\rho=(1-\rho)(I-\rho W)^{-1},
  \qquad B=S_\rho\,\xi .
\end{equation}

The restriction $0\leq\rho<1$ makes the Neumann expansion of $S_\rho$
a convex mixture of current and iterated peer averages.
Nonnegativity and unit row sums make each application of $W$ an
average, which is the property needed for Jensen's inequality below.
Because $W$ is directed, however, equal cross-sectional weights need
not be preserved by peer averaging.
Let $\pi$ instead be a strictly positive stationary distribution, so
$\pi^\top W=\pi^\top$.
Economically, $\pi_i$ measures firm $i$'s long-run influence under
repeated peer averaging.
Stationarity makes the aggregate mean invariant to $W$, while strict
positivity keeps every firm in the comparison and makes the weighted
quadratic norm nondegenerate.
For any exposure profile $z$, define

\begin{equation}\label{eq:cross-sectional-dispersion}
  V_\pi(z)=\sum_i\pi_i\norm{z_i-\bar z_\pi}^2,
  \qquad \bar z_\pi=\sum_i\pi_i z_i .
\end{equation}

Under these restrictions, peer adjustment cannot increase
cross-sectional dispersion.
It also cannot eliminate more than a coefficient-dependent share: at
least $\{(1-\rho)/(1+\rho)\}^2$ of stand-alone dispersion remains.
The finite, nonempty cross-section in the theorem keeps the weighted
mean and sums well defined.
The next result formalizes the bracket without requiring $W$ to be symmetric.

\begin{theorem}[Spatial attenuation]\label{thm:spatial-attenuation}
  Let $W$ be nonnegative and row stochastic, let $\pi$ be a strictly positive
  stationary distribution, and let $0\leq\rho<1$. For every finite nonempty
  cross-section $z$,
  \begin{equation}\label{eq:attenuation}
    \Big(\frac{1-\rho}{1+\rho}\Big)^{2}V_\pi(z)
    \;\leq\;
    V_\pi(S_\rho z)
    \;\leq\;
    V_\pi(z).
  \end{equation}
\end{theorem}
To establish the upper bound, let $f=z-\bar z_\pi\1$ and define
$\lVert f\rVert_\pi^2=V_\pi(z)$.
Stationarity makes $Wf$ centered whenever $f$ is, while nonnegative
unit rows permit Jensen's inequality:
\[
  \sum_i\pi_i\lVert(Wf)_i\rVert^2
  \leq\sum_{i,j}\pi_i W_{ij}\lVert f_j\rVert^2
  =\sum_j\pi_j\lVert f_j\rVert^2.
\]
The final equality uses stationarity, so $W$ contracts the centered
norm and hence so does every $W^k$.
Because $0\leq\rho<1$, the Neumann mixture
$S_\rho=(1-\rho)\sum_{k\geq0}\rho^k W^k$ is a convex combination of
these contractions.
This proves the upper bound.
For the lower bound, the centered profile $g=S_\rho f$ satisfies
$(1-\rho)f=(I-\rho W)g$.
The triangle inequality and the contraction $\lVert
Wg\rVert_\pi\leq\lVert g\rVert_\pi$ give
\[
  (1-\rho)\lVert f\rVert_\pi
  \leq\lVert g\rVert_\pi+\rho\lVert Wg\rVert_\pi
  \leq(1+\rho)\lVert g\rVert_\pi.
\]
Rearranging and squaring gives the lower bound in \eqref{eq:attenuation}.

At $\rho=0$, the two bounds coincide and peer adjustment leaves
dispersion unchanged.
For every admissible $\rho$, peer-adjusted exposures cannot be more
dispersed than stand-alone exposures, while the guaranteed retained
fraction $\{(1-\rho)/(1+\rho)\}^2$ declines with $\rho$.
Evaluating this fraction at the fitted working-model spatial-feedback
coefficient $\widehat\rho$ gives a descriptive plug-in lower
envelope, not the exact amount of attenuation or a calibration of
$t$, $L$, or $\tau_i$.
It has a structural interpretation as a retained-share bound only
under the return-projection conditions stated in the next section.
The bracket needs neither symmetry nor an additional spectral-gap
restriction, but it need not be sharp for the fitted operator.

The transfer and attenuation bounds can now be composed because they
measure dispersion in the same risk coordinates.
Set $Y_i=\Gamma^{1/2}B_i$, the risk-coordinate form of peer-adjusted
exposure, so linearity gives $Y=S_\rho Z$.
Then set the dispersion weights $q=\pi$.
This is not an arbitrary aggregation choice: the stationary weights
both preserve the mean under the directed operator and determine each
firm's contribution to the characteristic dispersion floor.
\begin{corollary}[Spatial-dispersion certificate]
  \label{cor:spatial-dispersion-certificate}
  Under the carrier and slack conditions stated above and the hypotheses of
  \Cref{thm:spatial-attenuation},
  \begin{equation}\label{eq:spatial-dispersion-certificate}
    \mathbb E\,V_\pi(Y)
    \geq
    \Big(\frac{1-\rho}{1+\rho}\Big)^2
    \left[L^{-1}\sqrt{\mathcal D_\pi(C_1,\dots,C_N)}-\tau_\pi\right]_+^2.
  \end{equation}
\end{corollary}
The complete chain is visible in one display.
The identity $V_\pi(Z)=\sum_{i<j}\pi_i\pi_j\lVert Z_i-Z_j\rVert^2$
links weighted variance to the multimarginal objective, and the
actual joint law of $Z$ is one feasible coupling.
Therefore
\[
  \begin{aligned}
    \mathbb E\,V_\pi(Y)
    &\geq
    \Big(\frac{1-\rho}{1+\rho}\Big)^2\mathbb E\,V_\pi(Z) \\
    &\geq
    \Big(\frac{1-\rho}{1+\rho}\Big)^2
    \mathcal D_\pi(P_1,\ldots,P_N) \\
    &\geq
    \Big(\frac{1-\rho}{1+\rho}\Big)^2
    \left[
      L^{-1}\sqrt{\mathcal D_\pi(C_1,\ldots,C_N)}-\tau_\pi
    \right]_+^2.
  \end{aligned}
\]
The first step applies attenuation realization by realization, the
second uses feasibility of the actual stand-alone exposure coupling,
and the third squares the transfer inequality.
This ordering shows how observed information heterogeneity becomes a
stand-alone exposure floor and then a peer-adjusted exposure floor
under the common adjustment coefficient $\rho$.
Information separation, carrier distortion, transmission noise, and
peer adjustment each enter once with a known direction.

The certificate is expressed in the latent risk coordinates of
peer-adjusted exposure, whereas the application observes scalar returns.
The next section states the return projection needed to preserve $W$
and the common adjustment coefficient $\rho$, then distinguishes that
structural parameter from the working-model adjustment index estimated by QMLE.
